% Options for packages loaded elsewhere
\PassOptionsToPackage{unicode}{hyperref}
\PassOptionsToPackage{hyphens}{url}
\PassOptionsToPackage{dvipsnames,svgnames,x11names}{xcolor}
\documentclass[
  11pt,
  a4paper,
]{article}
\usepackage{xcolor}
\usepackage[margin=24mm]{geometry}
\usepackage{amsmath,amssymb}
\setcounter{secnumdepth}{5}
\usepackage{iftex}
\ifPDFTeX
  \usepackage[T1]{fontenc}
  \usepackage[utf8]{inputenc}
  \usepackage{textcomp} % provide euro and other symbols
\else % if luatex or xetex
  \usepackage{unicode-math} % this also loads fontspec
  \defaultfontfeatures{Scale=MatchLowercase}
  \defaultfontfeatures[\rmfamily]{Ligatures=TeX,Scale=1}
\fi
\usepackage{lmodern}
\ifPDFTeX\else
  % xetex/luatex font selection
\fi
% Use upquote if available, for straight quotes in verbatim environments
\IfFileExists{upquote.sty}{\usepackage{upquote}}{}
\IfFileExists{microtype.sty}{% use microtype if available
  \usepackage[]{microtype}
  \UseMicrotypeSet[protrusion]{basicmath} % disable protrusion for tt fonts
}{}
\makeatletter
\@ifundefined{KOMAClassName}{% if non-KOMA class
  \IfFileExists{parskip.sty}{%
    \usepackage{parskip}
  }{% else
    \setlength{\parindent}{0pt}
    \setlength{\parskip}{6pt plus 2pt minus 1pt}}
}{% if KOMA class
  \KOMAoptions{parskip=half}}
\makeatother
\usepackage{longtable,booktabs,array}
\newcounter{none} % for unnumbered tables
\usepackage{calc} % for calculating minipage widths
% Correct order of tables after \paragraph or \subparagraph
\usepackage{etoolbox}
\makeatletter
\patchcmd\longtable{\par}{\if@noskipsec\mbox{}\fi\par}{}{}
\makeatother
% Allow footnotes in longtable head/foot
\IfFileExists{footnotehyper.sty}{\usepackage{footnotehyper}}{\usepackage{footnote}}
\makesavenoteenv{longtable}
\setlength{\emergencystretch}{3em} % prevent overfull lines
\providecommand{\tightlist}{%
  \setlength{\itemsep}{0pt}\setlength{\parskip}{0pt}}
% Shared style; used by generated, standalone LaTeX sources.
\usepackage{xcolor}
\usepackage{xurl}
\usepackage{fvextra}
\usepackage{etoolbox}
\usepackage{fancyhdr}
\usepackage{titlesec}
\definecolor{researchblue}{HTML}{173B58}
\AtBeginDocument{\hypersetup{linkcolor=researchblue,urlcolor=researchblue}}
\pagestyle{fancy}
\fancyhf{}
\fancyhead[L]{\small\sffamily SVG PATCH LAB}
\fancyhead[R]{\small\sffamily CVPR 2027 research package}
\fancyfoot[L]{\small Working documents; results and proposals distinguished}
\fancyfoot[R]{\thepage}
\setlength{\headheight}{15pt}
\setlength{\emergencystretch}{3em}
\titleformat{\section}{\Large\bfseries\color{researchblue}}{\thesection}{0.7em}{}
\titleformat{\subsection}{\large\bfseries\color{researchblue}}{\thesubsection}{0.7em}{}
\DefineVerbatimEnvironment{verbatim}{Verbatim}{breaklines=true,breakanywhere=true,fontsize=\small}
\AtBeginEnvironment{longtable}{\small\setlength{\tabcolsep}{4pt}}
\let\originaltableofcontents\tableofcontents
\renewcommand{\tableofcontents}{\originaltableofcontents\clearpage}
\usepackage{bookmark}
\IfFileExists{xurl.sty}{\usepackage{xurl}}{} % add URL line breaks if available
\urlstyle{same}
\hypersetup{
  pdftitle={Literature Survey and Reuse Decisions},
  pdfauthor={SVG Patch Lab - Research Working Package},
  colorlinks=true,
  linkcolor={Maroon},
  filecolor={Maroon},
  citecolor={Blue},
  urlcolor={blue},
  pdfcreator={LaTeX via pandoc}}

\title{Literature Survey and Reuse Decisions}
\author{SVG Patch Lab - Research Working Package}
\date{14 September 2026 \textbar{} CVPR 2027 target}

\begin{document}
\maketitle

{
\setcounter{tocdepth}{2}
\tableofcontents
}
Updated 13 September 2026. This supplements the earlier graphics survey
with the user's clarified focus: engineering problems solved using
physics and FEM.

\section{Conclusion}\label{conclusion}

\textbf{The broad idea already exists, including FEM agents, solver
verification, and fine-tuning on verified FEM programs.} A paper cannot
claim those combinations as new. The candidate extension is maintaining
numerical and semantic correctness in the final editable engineering
drawing, including after physical and presentation edits. This remains a
hypothesis to test against existing simulation and graphics systems; the
search does not establish unique novelty.

\section{Closest primary sources}\label{closest-primary-sources}

{\def\LTcaptype{none} % do not increment counter
\begin{longtable}[]{@{}
  >{\raggedright\arraybackslash}p{(\linewidth - 4\tabcolsep) * \real{0.3333}}
  >{\raggedright\arraybackslash}p{(\linewidth - 4\tabcolsep) * \real{0.3333}}
  >{\raggedright\arraybackslash}p{(\linewidth - 4\tabcolsep) * \real{0.3333}}@{}}
\toprule\noalign{}
\begin{minipage}[b]{\linewidth}\raggedright
Work
\end{minipage} & \begin{minipage}[b]{\linewidth}\raggedright
What overlaps
\end{minipage} & \begin{minipage}[b]{\linewidth}\raggedright
What to reuse or compare
\end{minipage} \\
\midrule\noalign{}
\endhead
\bottomrule\noalign{}
\endlastfoot
\textbf{MechAgents}, Ni \& Buehler, Extreme Mechanics Letters, 2024 &
LLM agents formulate, execute and correct FEM elasticity programs. & A
historical solver-agent baseline; published Colab notebooks require
older FEniCS/AutoGen dependencies. Do not call a modern reimplementation
an exact reproduction. \href{https://arxiv.org/abs/2311.08166}{Paper},
\href{https://github.com/lamm-mit/MechAgents}{author code}. \\
\textbf{MCP-SIM}, npj Artificial Intelligence, 2026 & Language-driven
simulation through tools and solver workflows. & Closest earlier
simulation orchestration comparison from the first survey;
runnable-release limitations are recorded there.
\href{https://www.nature.com/articles/s44387-025-00057-z}{Paper},
\href{https://github.com/KAIST-M4/MCP-SIM}{code}. \\
\textbf{FEABench}, Mudur et al., arXiv 2025; paper notes NeurIPS 2024
workshop acceptance & End-to-end natural-language engineering problems
solved through COMSOL APIs, with iterative tool feedback. & Reuse the
distinction between an executable solver call and a quantitatively
correct solution. COMSOL availability is a dependency for faithful
reproduction. \href{https://arxiv.org/abs/2504.06260}{Paper and linked
code}. \\
\textbf{FEM-Bench}, Mohammadzadeh et al., December 2025, revised May
2026 & FEM function generation and test generation against reference and
deliberately broken implementations. & \textbf{Selected for immediate
code reuse.} MIT-licensed tasks provide auditable numerical anchors. Its
33-task results are not directly comparable to our SVG editing scores.
\href{https://arxiv.org/abs/2512.20732}{Paper},
\href{https://github.com/elejeune11/FEM-bench}{repository}. \\
\textbf{ALL-FEM}, Deotale et al., CMAME 457:118985, 2026 & Fine-tuned
models generate FEniCS code in an agentic, feedback-driven workflow. The
authors report over 1,000 verified scripts, models from 3B to 120B, and
39 engineering benchmarks. & The most direct precedent for the proposed
training component. Training a small solver agent alone is not a new
contribution. Its reported code-level success is not a final-drawing
correctness metric. \href{https://arxiv.org/abs/2603.21011}{Paper},
\href{https://fenics-llm.github.io/}{author project and journal
citation}. \\
\textbf{PDEAgent-Bench}, May 2026 & 645 PDE instances across 11
families, with DOLFINx, Firedrake and deal.II tracks; numerical accuracy
and efficiency beyond mere execution. & Use its calibrated-reference
philosophy and PDE families when expanding beyond elementary anchors.
Keep benchmark test cases out of training.
\href{https://arxiv.org/abs/2605.09636}{Paper},
\href{https://github.com/YusanX/pde-agent-bench}{code}. \\
\textbf{OASiS / Open FEM Agent}, 2026 software & MCP access to multiple
FEM solvers, numerical verification, convergence studies and
visualization. & A current open-source orchestration baseline; its
scikit-fem backend matches our local solver ecosystem. This is software
prior art, not an assumed peer-reviewed performance result. We have
inspected its documentation, not reproduced its full agent evaluation.
\href{https://github.com/Hereon-InstituteMS/OASiS}{Repository and
citation}, \href{https://zenodo.org/records/20388035}{archived earlier
release}. \\
\textbf{VFEAgent}, May 2026, revised August 2026, preprint & Images and
text are converted into structured FEA specifications, then simulation
code with verification and debugging. & Directly relevant to CVPR
positioning: multimodal engineering-drawing-to-FEA workflows already
exist. Compare perception, physical specification and final artifact
separately. \href{https://arxiv.org/abs/2605.28978}{Paper}. \\
\textbf{ModSolAgent}, IEEE Transactions on Industrial Informatics, June
2026 & Abaqus script generation with retrieval and iterative
verification, plus distilled fine-tuning data. & Another direct
precedent for training engineering code models. Abaqus licensing and
released assets must be checked before reproduction.
\href{https://ieeexplore.ieee.org/document/11452232/}{Publisher}. \\
\end{longtable}
}

The prior \href{../../reports/PRIOR_ART.md}{graphics survey} remains
relevant for Penrose, scientific SVG generation and diagram editing. The
engineering papers establish the solver side; graphics papers establish
the editable-artifact side. Their combination still needs a precisely
defined, experimentally justified contribution.

\section{What has actually been built on top of prior
work}\label{what-has-actually-been-built-on-top-of-prior-work}

The vendored source snapshot \texttt{vendor/FEM-bench} is pinned to
commit \nolinkurl{d370aef5dbe023b0b86eb3548d24d11fd69328ee}; its MIT
license is retained. The new
\nolinkurl{scripts/fem_bench_svg_extension.py} directly imports the
author's unchanged \nolinkurl{FEM_1D_linear_elastic_CC0_H0_T0} reference
function.

The bridge passes both upstream reference tests, solves 12 generated
axial-bar cases, and performs 12 additional FEM solves after a 50\% load
increase. It writes displacements, reactions, provenance IDs and
editable SVG plots for both states. All 24 solves agree with the
analytic bar solution to below 2 \(\times\) 10\(^{-18}\) m in this run;
maximum relative force imbalance is below 4.2 \(\times\) 10\(^{-14}\).
These are elementary numerical checks, not evidence of general
engineering capability. A stale old tip displacement would be 33.3\% low
relative to the edited-load solution.

This is a working integration anchor. It does not reproduce the full
FEM-Bench leaderboard, train a model on FEM-Bench, or demonstrate a new
visual reasoning result. No upstream agent framework has been silently
relabeled as ours. Results are in \texttt{runs/fem-bench-svg-anchor/}.

Separately, \nolinkurl{scripts/robin_fem_reference.py} uses scikit-fem
to implement the weak form for the exact historical convective-plate
problem. It checks an analytic Robin anchor, discrete heat balance and
five mesh resolutions. This is an independent FEM implementation, not a
reproduction of MechAgents or ALL-FEM. See the
\href{../../reports/ASTRA_FEM_FAILURE_RECHECK.md}{Astra recheck}.

\section{Next contribution to test}\label{next-contribution-to-test}

Use a structured physical specification and solver-produced numerical
record as the source for an editable drawing. A physical change must
cause an actual new solve; a presentation change must preserve the
numerical claim and visible meaning. Measure errors after the final SVG
is rendered and edited, not only in a solver array or API response.

The main comparison should include direct Astra SVG generation, a
published solver-agent workflow, deterministic solver plotting,
unrestricted solver-informed SVG editing, and protected
numerical/semantic editing. Equalize available physical information and
count every attempted case. Test real geometry and material changes, not
only load scaling. Include human-written requests, multi-step edits,
quantity/unit/legend errors, and visual layout quality. Training is
justified only after a stronger task shows a gap that simple prompting
or a deterministic parser does not already close.

Prioritize well-posed Robin heat transfer, finite-domain elasticity with
explicit material/support/load assumptions, and electrostatics with a
stated dimensional and far-boundary model. A sharp-corner stress
singularity has no finite mesh-independent peak, so it is a
disclosure/identifiability test unless a fillet and load/support model
are specified. Reference uncertainty and singular regions must be
declared before model scoring.

For CVPR, a language-only code router is insufficient evidence. The
study needs an actual visual or graphics contribution and comparisons to
multimodal FEA and scientific diagram-editing work. The completed A100
pilot is useful infrastructure; it is not the main FEM experiment.

\clearpage

\section{Graphics and adjacent prior
art}\label{graphics-and-adjacent-prior-art}

{\def\LTcaptype{none} % do not increment counter
\begin{longtable}[]{@{}
  >{\raggedright\arraybackslash}p{(\linewidth - 4\tabcolsep) * \real{0.3333}}
  >{\raggedright\arraybackslash}p{(\linewidth - 4\tabcolsep) * \real{0.3333}}
  >{\raggedright\arraybackslash}p{(\linewidth - 4\tabcolsep) * \real{0.3333}}@{}}
\toprule\noalign{}
\begin{minipage}[b]{\linewidth}\raggedright
Work
\end{minipage} & \begin{minipage}[b]{\linewidth}\raggedright
Verified overlap
\end{minipage} & \begin{minipage}[b]{\linewidth}\raggedright
How this project could extend it
\end{minipage} \\
\midrule\noalign{}
\endhead
\bottomrule\noalign{}
\endlastfoot
\textbf{MechAgents}, Ni and Buehler, preprint 2023; Extreme Mechanics
Letters 2024 & Agents write, execute, and correct FEM code for
elasticity. Its examples explicitly include plotting the wrong stress
quantity and subsequently correcting it. & Reuse its elasticity tasks
and solver-agent baseline; automatically check that the displayed
quantity and contours agree with the solution, including after user
edits. \href{https://arxiv.org/abs/2311.08166}{Paper},
\href{https://github.com/lamm-mit/MechAgents}{author code and Colab
notebooks}. \\
\textbf{MCP-SIM}, Park, Moon and Ryu, npj Artificial Intelligence,
January 2026 & Natural-language specifications become simulations and
explanations through clarification, execution, and correction. Its
reported twelve-task success is specific to that evaluation. & Use its
workflow as a named simulation baseline. Add explicit solver-to-SVG
provenance and post-edit verification.
\href{https://www.nature.com/articles/s44387-025-00057-z}{Paper},
\href{https://github.com/KAIST-M4/MCP-SIM}{author repository}. \\
\textbf{SSVG-Bench / LOOP}, \emph{Towards the Generation of Structured
Scientific Vector Graphics with Large Language Models}, inspected ICLR
2026 submission & Scripted structural evaluation for plane geometry and
molecular diagrams, spanning TikZ, SVG and EPS; reasoning-oriented
prompting. The accessed manuscript is labeled under review; acceptance
and released code were not verified. & Extend structural checks to
numerical PDE fields and edit sequences; do not claim to invent
non-visual correctness metrics for scientific diagrams.
\href{https://openreview.net/pdf?id=aWypD2TAaC}{Manuscript}. \\
\textbf{Penrose}, SIGGRAPH 2020 & Separates mathematical meaning and
visual styling, uses constraints for layout, and exports SVG. & Borrow
the separation of semantics from layout. Physical contours must be
computed and bound to the solution; layout constraints alone do not
establish PDE accuracy.
\href{https://penrose.ink/media/Penrose_SIGGRAPH2020.pdf}{Author-hosted
paper}, \href{https://github.com/penrose/penrose}{code}. \\
\textbf{AutoFigure-Edit}, ACL system demonstrations 2026 & Generates
editable scientific illustrations with reference-guided styling and an
SVG editing interface. & Strong available graphics baseline: test
whether its edits preserve solver-derived geometry, field labels and
units. \href{https://aclanthology.org/2026.acl-demo.6/}{Published
paper}, \href{https://github.com/ResearAI/AutoFigure-Edit}{author
code}. \\
\textbf{SciDiagramEdit}, July 2026 preprint & Edits scientific vector
sources using instructions derived from paper revisions; improves an
editor through skill evolution over execution traces. & Closest
editing-learning comparison. Add numerical verification to the editing
objective. Its skill evolution is not the same as weight fine-tuning. No
official code release was verified in this audit.
\href{https://arxiv.org/abs/2607.15272}{Paper}. \\
\textbf{SVGEditBench V2}, February 2025 preprint & Instruction-based SVG
editing on emoji-derived examples, evaluated with visual, semantic and
geometric metrics. & Retain it as a general editing control, and add
physical-claim checks on engineering cases. It is not a physics
benchmark. \href{https://arxiv.org/abs/2502.19453}{Paper}. \\
\textbf{FMforME}, \emph{A Three-Layer Runtime Constraint Verification
Framework\ldots{}}, Applied Sciences 2026 & Validates structured CAD
specifications and feeds detected defects back to the LLM. Its stated
future work includes extending through FEA solving. & Compare
specification validity with actual solved-field fidelity. Do not copy
its domain-specific engineering rules as universal physical laws.
\href{https://www.mdpi.com/2076-3417/16/15/7396}{Article}. \\
\textbf{SciForma}, July 2026 preprint & Scientific methodology-diagram
generation with component, arrow and text evaluation and
multidimensional preference training. & Borrow the idea that all
required correctness dimensions must pass. Its released renderer is a
diffusion image model, so it is an adjacent visual baseline rather than
a native SVG/PDE implementation.
\href{https://arxiv.org/abs/2607.18091}{Paper},
\href{https://github.com/microsoft/SciForma}{author repository}. \\
\textbf{Text2PDE}, October 2024 preprint & Text-conditioned
latent-diffusion physics simulation. & Relevant if extending into
learned simulation, but much broader than the current need for
trustworthy SVG export and editing.
\href{https://arxiv.org/abs/2410.01153}{Paper}. \\
\end{longtable}
}

\section{Public-code feasibility}\label{public-code-feasibility}

The inspected MCP-SIM main tree has six agent modules, a README and an
MIT license. Several modules import \texttt{utils.setup\_logger}, but
\texttt{utils.py} is absent from that tree. A complete orchestration
entry point and dependency manifest were also absent. Its executor runs
generated Python and scans output for error patterns; that module does
not itself establish numerical correctness. The release needs
integration work before it can be treated as a reproducible baseline.
This observation concerns the public code, not a claim that the paper's
private experiments were invalid.
\href{https://github.com/KAIST-M4/MCP-SIM}{Repository},
\href{https://github.com/KAIST-M4/MCP-SIM/blob/main/simulation_executor_agent.py}{executor}.

MechAgents supplies Colab notebooks, but its README specifies older
FEniCS, AutoGen and OpenAI versions. Use it as a reproduction source
with a separately recorded environment, rather than silently mixing its
dependencies into this project.
\href{https://github.com/lamm-mit/MechAgents}{Repository}.

SciForma's README describes a 9B diffusion backbone and an eight-B200
training setup, with other distributed configurations in its tree. A
single Colab GPU would be a reduced adaptation experiment, not an
equivalent reproduction of that training setup.
\href{https://github.com/microsoft/SciForma}{Repository}.

\section{Paper archive}\label{paper-archive}

The catalog and bibliography are in \texttt{papers/}. Downloaded PDFs
retain their original authorship and publication notices. The download
manifest distinguishes bundled PDFs from source links and unavailable
downloads. This package does not bypass publisher access controls.

\end{document}
